\clearpage\thispagestyle{empty}
\pdfbookmark[0]{Quellen zu Instrumente (18. Oktober 1819)}{instrumentsoct1819-sources}
\begin{center}\textbf{Quellen zu Instrumente (18. Oktober 1819)}\end{center}
\noindent Diese KI-Übertragung ins heutige Deutsch folgt \textit{Instrumente} in den \href{https://gdz.sub.uni-goettingen.de/id/PPN236007467}{\textit{Werken}, Band~VI} (Göttingen, 1874), S.~422--424. Der Bericht erschien zuerst unter \textit{Göttingen} in den \textit{Göttingischen gelehrten Anzeigen}, Stück~167, 18. Oktober 1819. Die Werkausgabe nennt S.~1665--1669. Im \href{https://books.google.com/books?id=SggWAAAAYAAJ&pg=PA1665}{von Google digitalisierten Exemplar der New York Public Library} reicht der Bericht bis auf \href{https://books.google.com/books?id=SggWAAAAYAAJ&pg=PA1670}{S.~1670}. Dort stehen die Beobachtungen von 85 Pegasi und der Schlussabsatz über die Eigenbewegung des Sterns.
\par\smallskip\noindent Der Verweis am Anfang auf den \hyperlink{native-415}{früheren Bericht über den Meridiankreis von Repsold} steht in der Werkausgabe. Die Übertragung folgt der Schreibweise des wahrscheinlichen Fehlers in der Werkausgabe und behält die Ziffer 1 in der Bezeichnung des Sterns im Drachen und die fehlende Präposition in der Wendung über die Rolle und die Zimmerdecke bei. Der Bericht verwendet die erste Person Plural und spricht von Gauß in der dritten Person. KI-Übertragung und KI-Satz.
\clearpage
